\section*{第 \ref{ch:g4:symmetry} 章——对称}

\begin{solution}{exo:g4:symmetry:1}
正方形：是（$4$ 条\omterm{def:g4:symmetry:def}{对称轴}）。长方形：是（$2$ 条）。不等边三角形：否，沿任何直线
对折都不能重合。圆：是，过圆心的任意直线都是\omterm{def:g4:symmetry:def}{对称轴}。
\end{solution}

\begin{solution}{exo:g4:symmetry:2}
正方形：$4$ 条（两条中线和两条对角线）。长方形：$2$ 条（仅
两条中线）。等边三角形：$3$ 条（每条经过一个
\omterm{def:g4:geometry:polygon}{顶点}）。
\end{solution}

\begin{solution}{exo:g4:symmetry:3}
沿对角线对折后，两个三角形半块的形状相同，
但不能互相重合，其中一个“方向放反了”：描下的图形
落在原图旁边，而不是落在原图上。
因此对角线不是\omterm{def:g4:symmetry:def}{对称轴}。
\end{solution}

\begin{solution}{exo:g4:symmetry:4}
像 L 由 $2$ 个正方形组成，位于\omterm{def:g4:symmetry:def}{对称轴}的\emph{另一}侧，
方向左右颠倒（即镜中的 L），大小相同。
\end{solution}

\begin{solution}{exo:g4:symmetry:5}
补全后的图形在\omterm{def:g4:symmetry:def}{对称轴}下方画出旗帜及其上下颠倒的对称像，
旗杆与旗杆上下相对。
\end{solution}

\begin{solution}{exo:g4:symmetry:6}
$P$ 不动：沿\omterm{def:g4:symmetry:def}{对称轴}对折时，
位于\omterm{def:g4:symmetry:def}{对称轴}\emph{上}的每一点都保持原位，因此 $P$ 的对称像就是它自己。
\end{solution}

\begin{solution}{exo:g4:symmetry:7}
竖直\omterm{def:g4:symmetry:def}{对称轴}：A、H、O、T。水平\omterm{def:g4:symmetry:def}{对称轴}：B、E、H、O。
两者都有：H、O。两者都没有：F、N。
\end{solution}

\begin{solution}{exo:g4:symmetry:8}
是。这个三角形是等腰三角形，其\omterm{def:g4:symmetry:def}{对称轴}经过两条 $5$ cm 边之间的\omterm{def:g4:geometry:polygon}{顶点}，
并且经过 $3$ cm 底边的中点。对折时，
两条相等的边互换位置。
\end{solution}

\begin{solution}{exo:g4:symmetry:9}
任何恰有两条\omterm{def:g4:symmetry:def}{对称轴}的图形都可以，例如（非正方形的）长方形、
两臂长短不同的“十”字形（水平方向较长，竖直方向较短）、椭圆形。
两条\omterm{def:g4:symmetry:def}{对称轴}总相交于
图形的中心。
\end{solution}

\begin{solution}{exo:g4:symmetry:10}
完整的房子由左半边及其对称像组成：墙面宽度是半面墙的两倍，
上面是完整的三角形屋顶。
右边的每一条长度都与左边对应的长度相等，
对称会精确地复制长度。
\end{solution}

\begin{solution}{exo:g4:symmetry:11}
图中的画法说明了两种情况：等腰三角形沿经过\omterm{def:g4:geometry:polygon}{顶点}的直线对折，
两条相等的边 $\to$ \omterm{def:g4:symmetry:def}{对称轴}；等边三角形可以沿三个方向对折；
不等边三角形没有\omterm{def:g4:symmetry:def}{对称轴}（没有相等的边
$\to$ 没有\omterm{def:g4:symmetry:def}{对称轴}）。艾雅的两个结论都是正确的，
这一点将在\cref{ch:g6:symmetry}中证明。
\end{solution}
